\subsection{拓扑群的逼近}

设 G 是一个群，并设 \((\mathrm{H}_{\alpha})_{\alpha\in\mathbf{I}}\) 是 G 的正规子群的一个族，使得 \(H_{\alpha}\supset H_{\beta}\) 对一切 \(\alpha\leqslant\beta\) 成立。对每个 \(\alpha\in I\) ，令 \(G_{\alpha}=G/H_{\alpha}\) ；而对 \(\alpha\leqslant\beta\) ，令 \(f_{\alpha\beta}\) 为典范同态 \(G/H_{\beta}\to G/H_{\alpha}\) ，它因而把 \(H_{\beta}\) 在 G 中的一个陪集 T 映为陪集 \(TH_{\alpha}\) （\(H_{\alpha}\) 在 G 中的）。显然 \((G_{\alpha},f_{\alpha\beta})\) 是群的一个逆系统；\(\tilde{G}=\varprojlim G_{\alpha}\) 的元素是诸族 \((T_{\alpha})_{\alpha\in\mathbf{I}}\) ，其中 \(T_{\alpha}\) 是 \(H_{\alpha}\) 在 G 中的一个陪集（对每个 \(\alpha\) ），且 \(T_{\alpha}\supset T_{\beta}\) 对一切 \(\alpha\leqslant\beta\) 成立。映射 \(i:s\to(sH_{\alpha})\) （它把 G 映入 \(\tilde{G}\) ）是诸典范同态 \(G\to G/H_{\alpha}\) 的逆极限，因而是 G 到 \(\tilde{G}\) 中的一个同态，且元素 \((T_{\alpha})\in\tilde{G}\) 在 i 下的逆像等于 \(\bigcap_{\alpha\in I}T_{\alpha}\) 。于是 i 的核是 \(\bigcap_{\alpha\in I}H_{\alpha}\) ，而 i 的像由交非空的诸族 \((T_{\alpha})\in\tilde{G}\) 组成。

现在设 G 是一个拓扑群；若给每个 \(G_{\alpha} = G/H_{\alpha}\) 赋以商拓扑，则显然 \((\mathbf{G}_{\alpha}, f_{\alpha\beta})\) 是拓扑群的一个逆系统，且 \(i : G \to \tilde{G}\) 是一个连续同态。

命题 2. 设 G 是一个拓扑群，并设 \((\mathrm{H}_{\alpha})_{\alpha\in I}\) 是 G 的正规子群的一个族，使得 \(H_{\alpha}\supset H_{\beta}\) 对一切 \(\alpha\leqslant\beta\) 成立，且满足下面的条件：

(AP) 对每个 \(\alpha \in I\) ，\(H_{\alpha}\) 在 G 中闭，且 G 中单位元 \(e\) 的每个邻域都包含诸 \(H_{\alpha}\) 之一（换言之，诸 \(H_{\alpha}\) 组成的滤基收敛于 \(e\) ）。

那么，映射 \(i: \mathbf{G} \to \tilde{\mathbf{G}} = \varprojlim \mathbf{G} / \mathbf{H}_{\alpha}\) 是 \(\mathbf{G}\) 到 \(i(\mathbf{G})\) 上的一个严格态射；\(\tilde{\mathbf{G}}\) 是豪斯多夫的，且 \(i(\mathbf{G})\) 在 \(\tilde{\mathbf{G}}\) 中稠密；最后，\(i\) 的核是 \(\{e\}\) 在 \(\mathbf{G}\) 中的闭包。此外，若诸 \(\mathbf{H}_{\alpha}\) 中有一个是完备的，则 \(i\) 是满射的。

显然诸 \(G_{\alpha} = G/H_{\alpha}\) 是豪斯多夫的（§ 2，第 5 节，命题 13），因而 \(\tilde{G}\) 也是（因为它是 \(\prod_{\alpha \in I} G_{\alpha}\) 的一个子空间）。i 的核 H 是诸 \(H_{\alpha}\) 的交，因而是 G 的一个闭子群。由于 \(e\) 的每个邻域都包含某个 \(H_{\alpha}\) ，它包含 \(H\) ，因此{[}§ 3，第 1 节，公式 (1){]}\(H\) 是 \(\{e\}\) 的闭包。接下来我们证明 \(i(G)\) 在 \(\tilde{G}\) 中稠密。设 \(f_{\alpha}\) 是典范映射 \(\tilde{G} \to G_{\alpha}\) ，它是 \(\tilde{G}\) 上射影 \(pr_{\alpha}\) 的限制；\(\varphi_{\alpha} = f_{\alpha} \circ i\) 是典范映射 \(G \to G / H_{\alpha}\) 。若 \(U\) 是 \(\tilde{G}\) 的任一非空开集，则存在指标 \(\alpha \in I\) 与非空开集 \(U_{\alpha}\) （在 \(G_{\alpha}\) 中），使得 \(\overline{f}_{\alpha}^{-1}(U_{\alpha}) \subset U\) （第一章，§ 4，第 4 节，命题 9）；因而

\[
\overline {{i}} ^ {-1} (\mathbf {U}) \supset \overline {{\varphi}} _ {\alpha} ^ {-1} (\mathbf {U} _ {\alpha});
\]

但由于 \(\varphi_{\alpha}\) 是满射的，\(\overline{i}^{-1}(\mathbf{U})\) 非空，因而 \(i(\mathbf{G}) \cap \mathbf{U} \neq \emptyset\) 。为看出 \(i\) 是 \(\mathbf{G}\) 到 \(i(\mathbf{G})\) 上的一个严格态射，考虑一个邻域 \(\mathbf{V}\) （\(e\) 在 \(\mathbf{G}\) 中的）。存在一个邻域 \(\mathbf{W}\) （\(e\) 在 \(\mathbf{G}\) 中的），使得 \(\mathbf{W}^2 \subset \mathbf{V}\) ，又存在指标 \(\alpha \in \mathbf{I}\) ，使得 \(\mathbf{H}_{\alpha} \subset \mathbf{W}\) ；由此推出 \(\mathbf{V}\) 包含 \(\mathrm{WH}_{\alpha} = \overline{\varphi}_{\alpha}^{-1}(\varphi_{\alpha}(\mathbf{W})) = \overline{i}^{-1}(\overline{f}_{\alpha}^{-1}(\varphi_{\alpha}(\mathbf{W})))\) 。由于 \(\overline{f}_{\alpha}^{-1}(\varphi_{\alpha}(\mathbf{W}))\) 是 \(\tilde{\mathbf{G}}\) 中单位元的一个邻域，结论得证。

最后，假设存在指标 \(\gamma \in I\) ，使得 \(H_{\gamma}\) 是完备的。为证明 \(i\) 是满射的，只需证明每个族 \((T_{\alpha})_{\alpha \in I} \in \tilde{G}\) 有非空的交。由于 \(T_{\gamma}\) 是由 \(H_{\gamma}\) 经平移得到的，它是 \(G\) 的一个完备子空间（关于右一致结构与左一致结构两者）。此外，由于每个邻域 \(U\) （\(e\) 在 \(G\) 中的）都包含诸 \(H_{\alpha}\) 之一，相应的集合 \(T_{\alpha}\) 是 \(U_{d-}\) （或 \(U_{s-}\) ）小的，因而诸 \(T_{\alpha}\) （含于 \(T_{\gamma}\) 者）所成之集是一个柯西滤基；故它在 \(T_{\gamma}\) 中收敛，又由于诸 \(T_{\alpha}\) 在 \(G\) 中闭（因为它们是诸 \(H_{\alpha}\) 的平移），它们的交非空。

证毕。

推论 1. 若条件 (AP) 满足，且此外诸（豪斯多夫）群 \(\mathbf{G} / \mathbf{H}_{\alpha}\) 都是完备的，则群 \(\mathbf{G}\) 有一个可等同于 \(\tilde{\mathbf{G}}\) 的豪斯多夫完备化；且映射 \(i: \mathbf{G} \to \tilde{\mathbf{G}}\) 那时可等同于典范映射（§ 3，第 3 节，命题 5）。

因为此时 \(\tilde{\mathbf{G}}\) 是完备的（第 2 节），且命题 2 表明 \(i(\mathbf{G})\) 同构于由 \(\mathbf{G}\) 得到的豪斯多夫群；由于它在 \(\tilde{\mathbf{G}}\) 中稠密，推论得证（\(\S 3\) ，第 3 节，命题 5）。特别地：

推论 2. 设 G 是一个群，并设 \((\mathrm{H}_{\alpha})\) 是 G 的正规子群的一个族，它关于关系 \(\mathrm{H}_{\alpha} \supset \mathrm{H}_{\beta}\) 有向。若给 G 赋以使诸 \(\mathrm{H}_{\alpha}\) 构成单位元 \(e\) 的邻域基本系的群拓扑，则 G 的相伴豪斯多夫群同构于 \(\mathrm{G}_{1} = \mathrm{G} / \left( \bigcap_{\alpha} \mathrm{H}_{\alpha} \right)\) ；\(\mathrm{G}_{1}\) 有一个完备化 \(\hat{\mathrm{G}}_{1} = \tilde{\mathrm{G}}\) ；且典范映射 \(\mathrm{G}_{1} \to \hat{\mathrm{G}} = \varprojlim \mathrm{G} / \mathrm{H}_{\alpha}\) 可延拓为 \(\hat{\mathrm{G}} = \hat{\mathrm{G}}_{1}\) 到 \(\tilde{\mathrm{G}}\) 上的一个同构。

G 的子群 \(H_{\alpha}\) 是开的，因而也是闭的（\(\S\) 2，第 1 节，命题 4 的推论），且 \(G/H_{\alpha}\) 是离散的（\(\S\) 2，第 6 节，命题 18）；于是推论 1 的条件得到满足。

在本小节的其余部分，我们将假设 G 是豪斯多夫的，且 \((\mathbf{H}_{\alpha})\) 是 G 的紧正规子群的一个族，它关于关系 \(H_{\alpha} \supset H_{\beta}\) 有向并满足条件 (AP)；根据命题 2，映射

\[
i: \mathrm{G} \rightarrow \tilde {\mathrm{G}} = \varprojlim \mathrm{G} / \mathrm{H} _ {\alpha}
\]

于是是一个拓扑群的同构，它使我们可以把 G 与 \(\tilde{G}\) 等同起来。我们用 \(f_{\alpha}\) 表示典范映射 \(G \rightarrow G/H_{\alpha}\) 。

引理 1. 在命题 2 的假设下，若 E 是 G 的任一闭子集，我们有 \(\mathbf{E} = \bigcap_{\alpha} \mathrm{EH}_{\alpha}\) 。

因为 E 是诸集合 EV 的交，这里 V 取遍 e 的邻域滤子{[}§ 3，第 1 节，公式 (1){]}，且 e 的每个邻域都包含某个 \(H_{\alpha}\) ；又因 \(E \subset EH_{\alpha}\) ，结论得证。

命题 3. 假设 G 是豪斯多夫的，且诸 \(H_{\alpha}\) 是紧的并满足 (AP)。

\begin{enumerate}
\def\labelenumi{\alph{enumi})}
\item
  设 L 是 G 的一个闭子群；则对每个 \(\alpha\in I\) ，子群 \(L_{\alpha}=f_{\alpha}(L)\) （\(G_{\alpha}=G/H_{\alpha}\) 中的）是闭的，且 G 到 \(\varprojlim G_{\alpha}\) 上的同构 i 经限制给出 L 到 \(\varprojlim L_{\alpha}\) 上的一个同构。若还有 L 在 G 中正规，则 \(L_{\alpha}\) 在 \(G_{\alpha}\) 中正规（对每个 \(\alpha\in I\) ），且过渡到商，i 诱导出 G/L 到 \(\varprojlim G_{\alpha}/L_{\alpha}\) 上的一个同构。
\item
  反之，对每个 \(\alpha \in I\) ，设 \(L_{\alpha}\) 是 \(G_{\alpha}\) 的一个闭子群，使得 \(L_{\alpha} = f_{\alpha \beta}(L_{\beta})\) 对一切 \(\alpha \leqslant \beta\) 成立。则存在唯一的闭子群 \(L\) （\(G\) 中的），使得 \(L_{\alpha} = f_{\alpha}(L)\) 对每个 \(\alpha \in I\) 成立；而若此外 \(L_{\alpha}\) 在 \(G_{\alpha}\) 中正规（对每个 \(\alpha \in I\) ），则 \(L\) 在 \(G\) 中正规。
\item
  由于 \(\mathbf{H}_{\alpha}\) 是紧的，\(\mathrm{LH}_{\alpha}\) 在 \(\mathbf{G}\) 中闭（\(\S 4\) ，第 1 节，命题 1 的推论 1），因而 \(\mathbf{L}_{\alpha}\) 在 \(\mathbf{G}_{\alpha}\) 中闭。由于 \(i\) 把拓扑群 \(\mathbf{G}\) 与 \(\varprojlim \mathbf{G}_{\alpha}\) 等同起来，且 \(\varprojlim \mathbf{L}_{\alpha}\) 可与 \(\varprojlim \mathbf{G}_{\alpha}\) 的一个（拓扑）子群等同，\(i\) 把子群 \(\bigcap_{\alpha} \mathrm{LH}_{\alpha}\) （\(\mathbf{G}\) 的）与 \(\varprojlim \mathbf{L}_{\alpha}\) 等同起来，而为证明第一个断言，只需注意由引理 1 有 \(\mathbf{L} = \bigcap_{\alpha} \mathrm{LH}_{\alpha}\) 。另一方面，若 \(\mathbf{L}\) 正规，则对每个 \(\alpha \in \mathbf{I}\) ，映射 \(f_{\alpha}': \mathbf{G} / \mathbf{L} \to \mathbf{G}_{\alpha} / \mathbf{L}_{\alpha}\) （它由 \(f_{\alpha}\) 诱导）是一个满射的严格态射（\(\S 2\) ，第 8 节，注 3），其核是紧正规子群 \(\mathrm{H}_{\alpha} \mathrm{L} / \mathrm{L}\) （\(\mathrm{G} / \mathrm{L}\) 的），即紧子群 \(\mathrm{H}_{\alpha}\) （\(\mathbf{G}\) 的）的典范像。由于诸子群 \(\mathrm{H}_{\alpha} \mathrm{L} / \mathrm{L}\) （\(\mathrm{G} / \mathrm{L}\) 的）满足条件 (AP)（§ 2，第 8 节，命题 24），且 G/L 是豪斯多夫的，a) 的最后一个断言是命题 2 的推论。
\item
  设 \(f_{\alpha \beta}'\) 是 \(f_{\alpha \beta}\) 到 \(\mathbf{L}_{\beta} (\alpha \leqslant \beta)\) 上的限制。则 \((\mathbf{L}_{\alpha}, f_{\alpha \beta}')\) 是拓扑群的一个逆系统，其逆极限 \(\mathbf{L}\) 可与子群 \(\mathbf{G} \cap \prod_{\alpha} \mathbf{L}_{\alpha}\) （\(\mathbf{G}\) 的）等同。据假设，\(f_{\alpha \beta}'\) 是满射的，且其核是紧子群 \(f_{\beta}(\mathrm{H}_{\alpha}) \cap \mathrm{L}_{\beta}\) （\(\mathbf{L}_{\beta}\) 的）；因而（第一章，§ 9，第 6 节，命题 8 的推论 1）我们有 \(\mathbf{L}_{\alpha} = f_{\alpha}(\mathbf{L})\) ，对一切 \(\alpha \in \mathbf{I}\) 成立。若 \(\mathbf{L}'\) 是 \(\mathbf{G}\) 的另一个闭子群，使得 \(f_{\alpha}(\mathbf{L}') = \mathbf{L}_{\alpha}\) 对一切 \(\alpha \in \mathbf{I}\) 成立，则 \(\mathbf{L}'\mathbf{H}_{\alpha} = \overline{f}_{\alpha}^{-1}(\mathbf{L}_{\alpha})\) ，从而（引理 1）\(\mathbf{L}' = \bigcap_{\alpha} \mathbf{L}'\mathbf{H}_{\alpha} = \bigcap_{\alpha} \overline{f}_{\alpha}^{-1}(\mathbf{L}_{\alpha}) = \mathbf{L}\) 。最后，b) 的最后一个断言由公式 \(\mathbf{L} = \bigcap_{\alpha} \overline{f}_{\alpha}^{-1}(\mathbf{L}_{\alpha})\) 得出，因为诸 \(\overline{f}_{\alpha}^{-1}(\mathbf{L}_{\alpha})\) 现在是 \(\mathbf{G}\) 的正规子群。
\end{enumerate}

命题 4. 假设 G 是豪斯多夫的，诸 \(H_{\alpha}\) 是紧的，且 (AP) 满足。若 \(C_{\alpha}\) 是 \(G_{\alpha} = G/H_{\alpha}\) 的单位分量，则 G 的单位分量 C 可与 \(\varprojlim C_{\alpha}\) 等同，且我们有 \(f_{\alpha}(C) = C_{\alpha}\) 。本命题是下述引理的推论：

引理 2. 设 G 是一个豪斯多夫拓扑群，H 是 G 的一个紧正规子群，\(\varphi\) 是典范映射 \(\mathbf{G} \to \mathbf{G} / \mathbf{H}\) 。若 C 是 G 的单位分量，则 \(\varphi(\mathbf{C})\) 是 \(\mathbf{G}' = \mathbf{G} / \mathbf{H}\) 的单位分量。

一旦建立了这个引理，我们就将有 \(f_{\alpha}(\mathbf{C}) = \mathbf{C}_{\alpha}\) （对每个 \(\alpha \in \mathbf{I}\) ），且由于 \(\mathbf{C}\) 是 \(\mathbf{G}\) 的一个闭子群（§ 2，第 2 节，命题 7），只需应用命题 3a)。

为证明引理，首先注意，若 \(C'\) 是 \(e'\) 在 \(G'\) 中的连通分量，则 \(\varphi(C) \subset C'\) ，因为 \(\varphi(C)\) 是连通的。假设 \(\varphi(C) \neq C'\) 。由于 C 是 G 的一个闭正规子群（§2，第 2 节，命题 7），\(\varphi(C)\) 是 \(G'\) 的一个正规子群；若 \(\psi\) 是典范映射 \(G' \to G'/\varphi(C)\) ，则 \(\psi(C')\) 是连通的且不只由单位元组成；因而 \(G'/\varphi(C)\) 的单位分量不只由单位元组成。但 \(G'/\varphi(C)\) 同构于 \((G/H)/(HC/H)\) ，从而同构于 G/HC，于是也同构于 \((G/C)/(HC/C)\) （§2，第 7 节，命题 22 的推论以及命题 20）。现在，G/C 是豪斯多夫的且全不连通的（第一章，§11，第 5 节，命题 9），而 HC/C 作为 G 的紧正规子群 H 的典范像，是 G/C 的一个紧子群。因此只需在附加假设 G 全不连通、即 \(C = \{e\}\) 之下证明引理 2。

于是设 \(C' \neq \{e'\}\) ；把 G 换成它的子群 \(\varphi^{-1}(C')\) （它全不连通且包含 H），我们就可以假设 \(G'\) 是连通的且不止由一点组成。

设 \(\mathfrak{M}\) 是由闭子群 \(L\) （\(G\) 中的）组成的集合，这些子群满足 \(LH = G\) 。我们将证明 \(\mathfrak{M}\) 按关系 \(\supset\) 排序是归纳的。事实上，若 \(\mathfrak{X}\) 是 \(\mathfrak{M}\) 的一个线性有序子集，则对每个 \(x \in G\) ，诸集合 \(xH \cap L\) （\(L \in \mathfrak{X}\)）组成的集是一个由紧空间 \(xH\) 中的闭集组成的滤基；因而这些集合的交非空，这表明诸子群 \(L \in \mathfrak{X}\) 的交属于 \(\mathfrak{M}\) 。应用佐恩引理，我们因此看出 \(\mathfrak{M}\) 有一个极小元 \(L_0\) 。由于 \(H\) 是紧的，\(G/H = L_0H/H\) 同构于 \(L_0/(L_0 \cap H)\) （\(\S 4\) ，第 1 节，命题 1 的推论 3）；由于 \(L_0\) 全不连通且 \(L_0 \cap H\) 是紧的，我们看出可以把 \(G\) 换成 \(L_0\) ；换言之，我们可以附加假设不存在闭子群 \(L \neq G\) 满足 \(LH = G\) 。

现在设 F 是 G 中既开又闭的单位元的诸邻域的交，我们来证明 F 是 G 的一个闭子群。显然 F 在 G 中闭；因而只需证明 \(F^{-1}.F \subset F\) 。但若 \(x \in F\) 且 V 是 e 在 G 中的一个既开又闭的邻域，则 xV 也是；否则 e 将属于 xV 在 G 中的余集 W，而 W 又是既开又闭的，且我们将有 \(x \notin W\) ；从而 \(x \notin F\) ，与假设矛盾。由此推出 xF（诸集合 xV 的交，V 取遍 e 在 G 中的既开又闭的邻域）包含 F；换言之，\(x^{-1}F \subset F\) ，这就证明了我们的断言。由于 G 全不连通且 \(G \neq \{e\}\) ，我们有 \(F \neq G\) 。但若 V 是 e 在 G 中的一个既开又闭的邻域，则 VH 在 G 中也是既开又闭的（§ 4，第 1 节，命题 1 的推论 1），因而 \(\varphi(V)\) 在 G/H 中既开又闭，而这在假设之下蕴含 \(\varphi(V) = G/H\) 。我们将由此推出 FH = G，它将给出所要的矛盾并完成引理的证明。事实上，对每个 \(x \in G\) ，xH 与 e 的每个既开又闭的邻域 V 相交，因而也与这些邻域的交 F 相交，因为诸集合 \(Vx \cap H\) 组成一个由紧空间 xH 中的闭集组成的滤基。

证毕。

注. 若子群 \(H_{\alpha}\) 对某个 \(\alpha \in I\) 是紧的，则 \(H_{\beta}\) 对一切 \(\beta \geqslant \alpha\) 是紧的，因为它是 \(H_{\alpha}\) 的一个闭子群。由于所有 \(\beta \in I\) （满足 \(\beta \geqslant \alpha\) 的）组成的集合在 I 中共尾，在群 G 的研究中，假设诸 \(H_{\alpha}\) 之一是紧的与假设所有 \(H_{\alpha}\) 都是紧的，本质上有同样的结果。

